Randomised experiments on networks, such as platform experiments with social, marketplace or communication graphs, violate the no-interference assumption: a unit's outcome depends on the treatment of its neighbours. The usual target is the \emph{global average treatment effect} (GATE), the average outcome difference between treating everybody and treating nobody. The classical design-based answers (exposure mappings with Horvitz--Thompson or H\'ajek weights, graph-cluster randomisation) are unbiased under a correctly specified exposure model, but they need units to be fully exposed with non-negligible probability, which is rare in dense graphs, and they can have very large variance \citep{aronow2017estimating,ugander2013graph,eckles2017design}. A popular practical alternative is \emph{regression adjustment}: regress outcomes on one's own treatment and a summary of the neighbours' treatment, typically the treated fraction of neighbours, then predict the outcome with everybody treated and with nobody treated \citep{chin2019regression,han2023model}. It uses every unit, is cheap and has simple standard errors.

What is not well understood is \emph{how accurate it is, what drives its error, and when it breaks}. The prediction points lie far from where the design puts the data, so the estimator is an extrapolation whose accuracy depends on the design, the graph and the response shape; the cost of flexibility is variance. A practitioner who wants to trust or discard such an estimate needs to know which of these dominates in her setting.

We answer these questions in a controlled study with a known estimand. Outcomes are synthetic, so the GATE is exact, and are generated on four synthetic graph families and two public Facebook graphs; we vary the interference mechanism, the experimental design, graph size and density, interference strength, noise, the quality of the graph and the model class. Because every estimator sees the same assignments, differences are paired. Our contributions are:
\begin{itemize}
\item A reproducible benchmark of 72 main cells (6 graphs $\times$ 4 mechanisms $\times$ 3 designs) plus sweeps over eleven factors, with bootstrap uncertainty for every plotted number and a second seed set for replication (Section~\ref{sec:experiments}, Appendix).
\item An error decomposition into approximation (bias from the model class and the design's exposure support, measured with noise-free population fits) and estimation error, with three first-order analytic predictions (P1--P3) for bias and variance in the linear-exposure family that we check numerically on several hundred graph draws.
\item Findings with their scope: under a correct specification the estimator is unbiased and clearly better than the difference in means only for sufficiently large $n$ and sufficiently strong interference; its error is variance-dominated and grows with degree and with community-level confounding; with nonlinear exposure, wrong feature maps or graph errors it is biased by an amount that does not vanish with $n$, and flexible learners trade this bias for variance without dominating; cluster randomisation removes most of the advantage over the difference in means; and the standard cluster-robust intervals under-cover even when the point estimate is unbiased.
\item Negative results: a support diagnostic we expected to predict the bias works for one response family only; multi-$p$ designs and coverage trends under misspecification were not run.
\end{itemize}
All conclusions are about simulated outcomes on these graphs; we do not claim that they transfer to a given platform (Section~\ref{sec:discussion}).
